La electrificación del transporte público es una de las principales vías para reducir las emisiones y el ruido en las ciudades, y Santiago ha avanzado hasta operar una de las flotas de buses eléctricos más grandes fuera de China [citar fuente verificada]. Sin embargo, un bus eléctrico no se opera como uno diésel: su autonomía es limitada, su recarga toma tiempo y solo puede hacerse en electroterminales con pocos cargadores. Por esta razón, decidir qué bus hace cada viaje ya no basta, porque también hay que decidir cuándo y dónde carga cada bus, y ambas decisiones se condicionan entre sí.

Esa dependencia hace costosa una mala planificación, pues las horas de mayor demanda de viajes coinciden en parte con las de electricidad más cara y la recarga se concentra cuando los buses vuelven a los electroterminales. El resultado puede ser más buses de los necesarios, mayor gasto en energía y más kilómetros sin pasajeros, o incluso planes que no pueden ejecutarse porque la infraestructura no alcanza. Pero la misma complejidad abre oportunidades, ya que aprovechar mejor cada bus, por ejemplo haciéndolo atender viajes de más de una ruta, y repartir la recarga a lo largo del día puede reducir la flota, el gasto en energía y la presión sobre una infraestructura escasa y costosa de ampliar. Resolver bien este problema es, por lo tanto, una condición para que la electrificación sea sostenible desde el punto de vista económico y operacional.

Este informe continúa el Informe 1, que caracterizó el problema como un problema de asignación de vehículos eléctricos, conocido como E-VSP, y propuso una descomposición espaciotemporal con horizonte rodante. La revisión señaló que esa propuesta era demasiado compleja de implementar y pidió revisar en profundidad las estrategias de agrupamiento, por lo que esta entrega replantea la metodología y revisa distintas estrategias de agrupamiento de rutas.

En este contexto, el objetivo general es proponer y evaluar una metodología de planificación diaria de buses eléctricos que minimice el costo total de operación, y medir su desempeño frente a un caso base. Para lograrlo, se busca formular matemáticamente el problema, cuantificar el ahorro que genera permitir que un bus atienda viajes de rutas distintas, evaluar cómo la batería y los puestos de carga afectan el tamaño de la flota y el costo, estimar el valor de programar la carga frente a cargar solo cuando el bus lo necesita, e identificar los supuestos con mayor efecto sobre los resultados.

El entregable esperado es un plan operacional diario que resuelva las decisiones de asignación de buses a viajes, de traslados sin pasajeros y de recarga. Este plan se compone de las jornadas de cada bus, entendidas como la secuencia ordenada de expediciones y traslados que realiza entre su salida y su regreso al electroterminal, y de un programa de carga que especifica las ventanas de tiempo en que cada bus se recarga, el electroterminal utilizado y la energía repuesta. A esto se suma el desglose del costo total de la operación y los indicadores de desempeño calculados frente a un caso base.

